\chapter{Specification and Design}
\label{ch:design}

\section{Research Questions}

This project is guided by a comparative evaluation of
\glsxtrlong{icp}, \glsxtrlong{crc}, and \glsxtrlong{ltt}.
The primary research questions are:

\begin{itemize}
    \item How does \gls{icp} perform when adapted for a multi-label context,
    and what is the associated cost in terms of prediction set size and clinically relevant risks?

    \item What is the improvement in relevant risk control
    when using \gls{crc} over \gls{icp} for multi-label chest X-ray classification?

    \item How is the effectiveness of \gls{crc} impacted
    when applied to more nuanced risk functions
    where missing more critical pathologies incurs a higher penalty?

    \item How do CRC and LTT compare in their ability to control a clinically critical,
    monotone risk like the False Negative Rate (FNR)?
    Is the high-probability guarantee of LTT achievable without a significant loss of utility
    (i.e., overly large prediction sets) compared to CRC's guarantee in expectation?

    \item Leveraging the unique capability of the LTT framework,
    how effectively can the \glsxtrlong{fdr} (a non-monotone risk) be controlled
    for multi-label \gls{cxr} classification?

    \item How does the effect of CRC and LTT vary across different model architectures?


\end{itemize}

\section{Experimental Design}

The experimental design for this project will involve the following components:

\subsection{Datasets}

Two large-scale, publicly available chest X-ray datasets were chosen
to ensure the findings are robust and generalizable:
\begin{itemize}
    \item \textbf{CheXpert}~\cite{irvin2019chexpert}:
    A large dataset of chest X-rays with multi-label annotations for 14 different pathologies.
    It includes over 200,000 radiographs and is widely used in the literature.

    \item \textbf{NIH ChestX-ray14}~\cite{wang2017chestxray}:
    Another large dataset containing over 100,000 frontal-view chest X-ray images
    with 14 pathologies. It is also a standard benchmark in the field.

    \item \textbf{PadChest}~\cite{Bustos2020padchest}:
    A third dataset containing over 160,000 chest X-rays
    with multi-label annotations for 174 different pathologies.
\end{itemize}

\subsection{Models}

These models are widely used in the literature and provide a good baseline for comparison.

\begin{itemize}
    \item \emph{ResNet50}~\cite{he2015resnet}:
    A 50-layer Residual Network that was selected for being
    a foundational architecture in modern deep learning.
    Its use of residual connections allows for the effective training of very deep networks
    by mitigating the vanishing gradient problem.
    It serves as a strong and widely used baseline for classification tasks.

    \item \emph{DenseNet121}~\cite{huang2017densenet}:
    A 121-layer Densely Connected Convolutional Network that was chosen
    to represent an alternative architectural philosophy.
    DenseNet connects each layer to every other layer in a feed-forward fashion,
    promoting feature reuse, improving feature propagation, and often achieving
    high performance with fewer parameters than comparable ResNet models.

    \item \emph{Official CheXpert Model}~\cite{irvin2019chexpert}:
    A DenseNet inspired official baseline model provided by the dataset's authors was also included.
    This model is specifically pre-trained and configured for the CheXpert task.
\end{itemize}

\subsection{Metrics}

To evaluate the performance of each framework, the following metrics were collected:
\begin{itemize}
    \item Average Set Size:
    The average size of the prediction sets generated by each method. Smaller sets indicate
    higher confidence in predictions.

    \item Coverage (for \gls{cp}):
    The proportion of true labels that are included in the prediction sets.
    This metric assesses the ability of the method to capture all relevant labels.

    \item Risk Control (for \gls{crc} and \gls{ltt}):
    Calculated as the empirical loss for each method.
    For CRC, this is the average risk expected to be at or below target \(\alpha\).
    For LTT, this is the risk expected to be at or below the target \(\alpha\) with high probability.
\end{itemize}

Different loss functions and risk definitions (e.g., class-specific FNRs) will be explored
to study how they influence the behaviour of the methods and the generated prediction sets on this task.

The \texttt{torchxrayvision}~\cite{Cohen2022xrv} library will be used
to load the datasets and preprocess the images.
This library provides a convenient interface for loading and transforming chest X-ray images,
including resizing, normalization, and data augmentation.
The models will run under \texttt{pytorch},
a popular deep learning framework for machine learning and computer vision tasks.
